\section{Proofs for Section~\texorpdfstring{\ref{sec:heavy}}{4}}\label{app:heavy}

We use the notation of Section~\ref{sec:heavy}, with $\tau_0=0$. For a first
run $D$ let $\tau^D_m=D+\tau'_{m-1}$, where $(\tau'_m)$ is an independent copy
of $(\tau_m)$, and $\tau^D_0=0$. Then $u^D_m(k)=\Prob(\tau^D_m=k)$ and
$q^D_m(k)=\Prob(\tau^D_{m-1}<k\le\tau^D_m)$. Conditioning on the last run
gives, for $m,k\ge1$,
\begin{equation}\label{eq:C-conv}
\begin{gathered}
q_m(k)=\sum_{j\ge1}\bar F(j)\,u_{m-1}(k-j),\qquad
q^D_{m+1}(k)=\sum_{j\ge1}\bar F(j)\,u^D_m(k-j),\\
u^D_m(k)=\sum_{d\ge1}\Prob(D=d)\,u_{m-1}(k-d).
\end{gathered}
\end{equation}
Since $\tau_m$ and $\tau^D_m$ increase strictly, $\sum_mq_m(k)=\sum_mq^D_m(k)=1$,
$\sum_mu^D_m(k)\le1$ and $\sum_{m\le m'}q_m(k)=\Prob(\tau_{m'}\ge k)$. We put
$z_{m,k}=(k-m\mu)/a_m$ if $\alpha>1$ and $z_{m,k}=k/a_m$ if $\alpha<1$, and
$\delta_m=\sup_k\lvert a_mu_m(k)-g(z_{m,k})\rvert$. We write $\omega$ for the
modulus of continuity of $g$ and $\bar g(x)=\sup_{\lvert y\rvert\ge x}g(y)$.

\subsection{Proof of Lemma~\texorpdfstring{\ref{lem:stable}}{4.3}}

(i) Summation by parts gives $\E(1-x^L)=(1-x)\sum_{k\ge1}\bar F(k)x^{k-1}$,
for complex $\lvert x\rvert\le1$ if $\alpha>1$ (the sums converge absolutely
since $\sum_kkp_k=\mu$) and for $0\le x<1$ if $\alpha<1$ (the terms are
positive). With $x=e^{-u}$ this reads
\begin{equation}\label{eq:C-polylog}
1-\E e^{-uL}=(e^u-1)\,\mathrm{Li}_\alpha(e^{-u}),\qquad
\mathrm{Li}_\alpha(z)=\sum_{k\ge1}k^{-\alpha}z^k .
\end{equation}
For non-integer $\alpha$, $\operatorname{Re}u\ge0$ and $0<\lvert u\rvert<2\pi$,
\cite[25.12.12]{dlmf} gives
$\mathrm{Li}_\alpha(e^{-u})=\Gamma(1-\alpha)u^{\alpha-1}+\sum_{k\ge0}\zeta(\alpha-k)(-u)^k/k!$
with the principal branch. For $1<\alpha<2$ this gives
$1-\E e^{-uL}=\mu u+\Gamma(1-\alpha)u^\alpha+O(\lvert u\rvert^2)$, and
$u=-i\theta$ gives (i), since $K_\alpha=-\Gamma(1-\alpha)$.

(ii) With $t=\theta/a_m$, part (i) gives
$e^{-i\mu t}\E e^{itL}=1+K_\alpha(-i\theta)^\alpha/m+O(\theta^2m^{-2/\alpha})$.
Since $m^{1-2/\alpha}\to0$, the characteristic function of $(\tau_m-m\mu)/a_m$
tends to $\exp(K_\alpha(-i\theta)^\alpha)$, whose modulus
$\exp(K_\alpha\cos(\frac{\pi\alpha}2)\lvert\theta\rvert^\alpha)$ is
integrable. L\'evy's continuity theorem, Fourier inversion and the
Riemann--Lebesgue lemma give $G$ and the stated properties of $g$. For real
$u>0$, part (i) reads $\E e^{-uL}=1-\mu u+K_\alpha u^\alpha+O(u^2)$, and the
same computation gives $\E e^{-s(\tau_m-m\mu)/a_m}\to e^{K_\alpha s^\alpha}$
for $s\ge0$. Applied with $2s$ in place of $s$, this shows that these
variables are bounded in $L^2$. So they are uniformly integrable, and
$\E e^{-sG}=e^{K_\alpha s^\alpha}$.

(iii) Every $p_k$ is positive, so $L$ has maximal span 1. By (ii) and (v) the
normalized sums converge to stable laws, and Gnedenko's local limit theorem
\cite[Theorem~4.2.1]{ibragimovlinnik1971} gives $\delta_m\to0$. Then
$a_mu_m(k)\le\lVert g\rVert_\infty+1$ for large $m$, and $a_mu_m(k)\le a_m$
for the other $m$, so $c_0<\infty$.

(iv) If $G'$ is an independent copy of $G$, then
$\lvert\E e^{i\theta G}\rvert^2=e^{-\sigma_\alpha\lvert\theta\rvert^\alpha}$
is the characteristic function of $G-G'$. By Plancherel's theorem,
$I_\alpha=\frac1{2\pi}\int e^{-\sigma_\alpha\lvert\theta\rvert^\alpha}d\theta
=\Gamma(1+1/\alpha)/(\pi\sigma_\alpha^{1/\alpha})$.

(v) Let $0<\alpha<1$. For real $u\downarrow0$, \eqref{eq:C-polylog} gives
$1-\E e^{-uL}=\Gamma(1-\alpha)u^\alpha+O(u)$, with $\Gamma(1-\alpha)>0$. As
$1/\alpha>1$, $m(1-\E e^{-sL/a_m})\to\Gamma(1-\alpha)s^\alpha$, so
$\E e^{-s\tau_m/a_m}\to e^{-\Gamma(1-\alpha)s^\alpha}$, and the continuity
theorem for Laplace transforms gives $\tau_m/a_m\Rightarrow G$. Both sides of
$\E e^{-sG}=e^{-\Gamma(1-\alpha)s^\alpha}$ are continuous on
$\operatorname{Re}s\ge0$ and analytic on $\operatorname{Re}s>0$, so the
identity holds at $s=-i\theta$. Hence
$\lvert\E e^{i\theta G}\rvert=\exp(-\Gamma(1-\alpha)\cos(\frac{\pi\alpha}2)\lvert\theta\rvert^\alpha)$
is integrable, and $g$ exists as in (ii). Also $g=0$ on $(-\infty,0]$ since
$G\ge0$. For the tail, if $\E e^{-sG}=e^{-s^\alpha}$ then Pollard's series
\cite{pollard1946} writes the density as a series in the powers
$y^{-\alpha k-1}$, $k\ge1$, that converges for $y>0$, with first term
$\frac1\pi\Gamma(1+\alpha)\sin(\pi\alpha)\,y^{-1-\alpha}$. Scaling by
$\Gamma(1-\alpha)^{1/\alpha}$ and the reflection formula
$\Gamma(1+\alpha)\Gamma(1-\alpha)=\pi\alpha/\sin(\pi\alpha)$ give
$g(y)=\alpha y^{-1-\alpha}+O(y^{-1-2\alpha})$ as $y\to\infty$.
\qed

\subsection{Two lemmas on renewal sums}

A sum of $m$ runs that is much larger than its typical size is most likely
made by a single long run. The first lemma turns this into a local bound.

\begin{lemma}\label{lem:C-bigjump}
There is a constant $C_B=C_B(\alpha)$ such that
$u_m(k)\le C_B\,m\,k^{-\alpha-1}$ for all $m\ge1$ in each of the cases
\textup{(a)} $1<\alpha<2$ and $k\ge2m\mu$, \textup{(b)} $\alpha\ge2$ and
$k\ge2m\mu$, \textup{(c)} $0<\alpha<1$ and $k\ge a_m$.
\end{lemma}

\begin{proof}
Here $C$ depends only on $\alpha$ and may change from line to line. Since
$\tau_m\ge m$ we may assume $m\le k$. For $r>2$, chosen below, put $y=k/r$.

\emph{One big jump.} By the mean value theorem
$p_\ell\le\alpha\ell^{-\alpha-1}$. So for each $i$,
$\Prob(\tau_m=k,\,L_i\ge y)=\sum_{\ell\ge y}p_\ell\,\Prob(\tau_m-L_i=k-\ell)
\le\alpha r^{\alpha+1}k^{-\alpha-1}$, and the $m$ indices give at most
$\alpha r^{\alpha+1}mk^{-\alpha-1}$.

\emph{No big jump.} We may take $y\ge1$, since otherwise $\max_iL_i<y$ is
impossible. Let $I$ be a set of indices, $\tau_I=\sum_{i\in I}L_i$, $x_0>0$
and $\lambda=s/y$ with $s\ge0$. Markov's inequality and
$e^x\le1+x+\frac{x^2}2e^x$ for $x\ge0$ give
\begin{equation}\label{eq:C-chernoff}
\Prob\bigl(\tau_I\ge x_0,\ \max_{i\in I}L_i<y\bigr)
\le\exp\Bigl(-\frac{sx_0}y+\lvert I\rvert\Bigl(\frac sy\,\E[L;L<y]
+\frac{s^2e^s}{2y^2}\,\E[L^2;L<y]\Bigr)\Bigr).
\end{equation}
In case (a) we take all $m$ indices and $x_0=k$, and use $\E[L;L<y]\le\mu$,
$\E[L^2;L<y]\le\alpha\sum_{\ell<y}\ell^{1-\alpha}\le Cy^{2-\alpha}$ and
$k-m\mu\ge k/2$. The exponent is then at most
$-sr/2+Cr^\alpha s^2e^sm/k^\alpha$. With $s=\frac12\log(k^\alpha/m)$, which
is positive since $m\le k/2<k^\alpha$, we have $e^sm/k^\alpha=e^{-s}$ and
$s^2e^{-s}\le1$. So the bound is $C(m/k^\alpha)^{r/4}$. With
$r=4+4/(\alpha-1)$ and $m\le k$ this is at most
$C(m/k^\alpha)(k^{1-\alpha})^{1/(\alpha-1)}=Cmk^{-\alpha-1}$.

In case (b) we argue as in case (a). For $\alpha>2$ we use
$\E[L^2;L<y]\le\E L^2<\infty$,
$s=\frac12\log(k^2/m)$ and $r=4\alpha$, which give
$C(m/k^2)^\alpha\le Cmk^{-\alpha-1}$. For $\alpha=2$ we use
$\E[L^2;L<y]\le2\sum_{\ell<y}\ell^{-1}\le C\log(ek)$, $r=20$ and
$s=\frac12\log\frac{k^2}{m\log(ek)}$, which is positive since
$m\log(ek)\le k\log(ek)/(2\mu)<k^2$. This gives
$C\bigl(m\log(ek)/k^2\bigr)^5\le Cmk^{-6}\log^5(ek)\le Cmk^{-3}$.

In case (c) let $t=k^\alpha/m\ge1$, so that $k=t^{1/\alpha}a_m$, and
$r=4+4/\alpha$. For $m=1$ the claim is $p_k\le\alpha k^{-\alpha-1}$. For
$m\ge2$ split the indices into $I_1=\{1,\dots,\lfloor m/2\rfloor\}$ and
$I_2=\{\lfloor m/2\rfloor+1,\dots,m\}$. On $\{\tau_m=k\}$ one of
$\tau_{I_1},\tau_{I_2}$ is at least $k/2$, so by independence
\[
\Prob\bigl(\tau_m=k,\ \max_iL_i<y\bigr)\le\sum_{i=1,2}
\Prob\bigl(\tau_{I_i}\ge k/2,\ \max_{I_i}L<y\bigr)\,\sup_j\Prob(\tau_{I_{3-i}}=j).
\]
Each half has at least $\lfloor m/2\rfloor\ge m/3$ indices, so by
Lemma~\ref{lem:stable}(iii) the supremum is at most
$c_0/a_{\lfloor m/2\rfloor}\le3^{1/\alpha}c_0/a_m$. In \eqref{eq:C-chernoff}
with $x_0=k/2$ we use $\E[L;L<y]\le Cy^{1-\alpha}$,
$\E[L^2;L<y]\le Cy^{2-\alpha}$, $\lvert I_i\rvert\le m$ and
$my^{-\alpha}=r^\alpha/t$. The exponent is at most $-sr/2+C(s+s^2e^s)/t$,
and $s=\frac12\log t$ makes the second term bounded. So the first factor is at
most $Ct^{-r/4}$, and $t^{-r/4}/a_m=t^{-1-1/\alpha}/a_m=mk^{-\alpha-1}$.
\end{proof}

\begin{lemma}\label{lem:C-firstpassage}
Let $1<\alpha<2$ and let $D$ be any a.s.\ finite first-run law. Then
\[
\sup_k\Bigl\lvert\frac{a_m}\mu q_m(k)-g(z_{m,k})\Bigr\rvert,\qquad
\sup_k\bigl\lvert a_mu^D_m(k)-g(z_{m,k})\bigr\rvert,\qquad
\sup_k\Bigl\lvert\frac{a_m}\mu q^D_{m+1}(k)-g(z_{m,k})\Bigr\rvert
\]
tend to $0$ as $m\to\infty$. For $m\ge2$, $\sup_kq_m(k)\le\mu c_0/a_{m-1}$
and $\sup_ku^D_m(k)\le c_0/a_{m-1}$.
\end{lemma}

\begin{proof}
By \eqref{eq:C-conv}, $q_m$ and $u^D_m$ are mixtures of shifts of $u_{m-1}$
with total weights $\mu$ and $1$. This gives the bounds for $m\ge2$.

Fix $J\ge1$ and $Z>0$, put $T_J=\sum_{j>J}\bar F(j)$ and write $z=z_{m,k}$.
In the sum for $q_m(k)$ the terms with $j>J$ add up to at most
$c_0T_J/a_{m-1}$. For $j\le J$ we have
$\lvert a_{m-1}u_{m-1}(k-j)-g(z_{m-1,k-j})\rvert\le\delta_{m-1}$ and
$z_{m-1,k-j}=z\,a_m/a_{m-1}+(\mu-j)/a_{m-1}$. If $\lvert z\rvert\le Z$, then
$\lvert z_{m-1,k-j}-z\rvert\le\eta_m=Z(a_m/a_{m-1}-1)+(J+\mu)/a_{m-1}$. If
$\lvert z\rvert>Z$ and $(J+\mu)/a_{m-1}\le Z/2$, then $g(z)$ and
$g(z_{m-1,k-j})$ both lie in $[0,\bar g(Z/2)]$. Since
$\mu-\sum_{j\le J}\bar F(j)=T_J$, we get, uniformly in $k$ for large $m$,
\[
\bigl\lvert a_{m-1}q_m(k)-\mu g(z)\bigr\rvert
\le\mu\bigl(\delta_{m-1}+\omega(\eta_m)+\bar g(Z/2)\bigr)
+(c_0+\lVert g\rVert_\infty)T_J .
\]
We let $m\to\infty$, then $J,Z\to\infty$, and use $a_m/a_{m-1}\to1$. This
gives the first limit. The second follows in the same way from the sum for
$u^D_m$, with the weights $\Prob(D=d)$ and with $\Prob(D>J)$ in place of
$T_J$. For the third we apply the argument to the sum for $q^D_{m+1}$, with
$a_mu^D_m$ in place of $a_{m-1}u_{m-1}$. Then the second limit replaces
$\delta_{m-1}\to0$, the shift is $z_{m,k-j}-z_{m,k}=-j/a_m$, and the tail is
at most $c_0T_J/a_{m-1}$.
\end{proof}

\subsection{Proof of Theorem~\texorpdfstring{\ref{thm:center}}{4.4}}

We write the proof so that it also gives Theorem~\ref{thm:alpha2}. It uses
only the following facts. First, $a_m=a(m)$ for large $m$, where $a$ is
increasing and regularly varying with index in $(0,1)$. Second,
$(\tau_m-m\mu)/a_m\Rightarrow G$, where $G$ has a bounded, uniformly
continuous density $g$ that vanishes at $\pm\infty$, and $\delta_m\to0$ and
$c_0<\infty$. Third, Lemma~\ref{lem:C-firstpassage} holds and
$u_m(k)\le C_Bmk^{-\alpha-1}$ for $k\ge2m\mu$. For $1<\alpha<2$ these hold with
$a(x)=x^{1/\alpha}$, by Lemma~\ref{lem:stable} and
Lemma~\ref{lem:C-bigjump}(a). Put $b_h=a(h/\mu)$, so that
$b_h=(h/\mu)^{1/\alpha}$ for $1<\alpha<2$. We show that
$C_n\sim2b_h^{-1}\int g^2$. For $1<\alpha<2$ this is
$2I_\alpha(\mu/h)^{1/\alpha}$, and the second form in
Theorem~\ref{thm:center} follows from
$\sigma_\alpha=2K_\alpha\lvert\cos\frac{\pi\alpha}2\rvert$ and $h=n/2$.

By regular variation, $b_h=o(h)$, $a(h/(8\mu))\ge b_h/9$ and
$a(h/(2\mu))\ge b_h/3$ for large $h$, and $a_m/b_h\to1$ uniformly in $m$
with $\lvert m\mu-h\rvert\le2Mb_h$, for each fixed $M$. By
Proposition~\ref{prop:tworenewal},
$C_n=\sum_{m\ge1}\bigl[u_m(h)q^D_{m+1}(h)+u^D_m(h)q_m(h)\bigr]$. Fix $M>1$, put
$m_\pm=(h\pm Mb_h)/\mu$, and let $\Sigma_1,\dots,\Sigma_4$ be the parts of
this sum over the regions
\[
\mathcal R_1=\{m\le\tfrac h{4\mu}\},\quad
\mathcal R_2=\{\tfrac h{4\mu}<m<m_-\},\quad
\mathcal R_3=\{m_-\le m\le m_+\},\quad
\mathcal R_4=\{m>m_+\}.
\]

\emph{Region $\mathcal R_1$.} Here $h\ge4m\mu$, so
$u_m(h)\le C_Bmh^{-\alpha-1}\le C_Bh^{-\alpha}/(4\mu)$. Also
$q_1(h)=h^{-\alpha}$. For $m\ge2$ we split the sum for $q_m(h)$ at $j=h/2$.
The terms with $j\ge h/2$ add up to at most $(h/2)^{-\alpha}$. For $j<h/2$ we
have $h-j>h/2\ge2\mu(m-1)$, so the other terms add up to at most
$\mu C_B(m-1)(h/2)^{-\alpha-1}\le2^{\alpha-1}C_Bh^{-\alpha}$. Since
$\sum_mq^D_{m+1}(h)\le1$ and $\sum_mu^D_m(h)\le1$, we get
$\Sigma_1=O(h^{-\alpha})$. This is $o(1/b_h)$, since the index of $a$ is
less than $1<\alpha$.

\emph{Region $\mathcal R_2$.} For large $h$ we have $m-1\ge h/(8\mu)$ here, so
$a_m\ge a_{m-1}\ge b_h/9$, and $u_m(h),u^D_m(h)\le9c_0/b_h$ by
Lemma~\ref{lem:C-firstpassage}. Let $m'=\lfloor m_-\rfloor$. Then
\[
\sum_{m\in\mathcal R_2}q_m(h)\le\Prob(\tau_{m'}\ge h),\qquad
\sum_{m\in\mathcal R_2}q^D_{m+1}(h)\le\Prob(D+\tau'_{m'}\ge h)
\le\Prob(D\ge b_h)+\Prob(\tau_{m'}\ge h-b_h).
\]
Since $m'\mu\le h-Mb_h$ and $a_{m'}\le b_h$, the event
$\{\tau_{m'}\ge h-b_h\}$ implies $(\tau_{m'}-m'\mu)/a_{m'}\ge M-1$. As $D$ is
a.s.\ finite and $(\tau_{m'}-m'\mu)/a_{m'}\Rightarrow G$, we get
$\limsup_hb_h\Sigma_2\le18c_0\,\Prob(G\ge M-1)$.

\emph{Region $\mathcal R_4$.} For large $h$ we have $m-1\ge h/(2\mu)$ here, so
$u_m(h),u^D_m(h)\le3c_0/b_h$. Let $k=\lfloor m_+\rfloor$. Then
$\sum_{m\in\mathcal R_4}q_m(h)=\Prob(\tau_k<h)$, and since $D\ge1$,
$\sum_{m\in\mathcal R_4}q^D_{m+1}(h)=\Prob(D+\tau'_k<h)\le\Prob(\tau_k<h)$.
Now $(h-k\mu)/a_k\le(\mu-Mb_h)/a_k\to-M$, and the law of $G$ is continuous,
so $\limsup_hb_h\Sigma_4\le6c_0\,\Prob(G\le-M)$.

\emph{Region $\mathcal R_3$.} Here $m\to\infty$ and $a_m=b_h(1+o(1))$
uniformly. By Lemma~\ref{lem:C-firstpassage} and $\delta_m\to0$, uniformly in
$m\in\mathcal R_3$,
\[
u_m(h)q^D_{m+1}(h)+u^D_m(h)q_m(h)=\frac{2\mu}{a_m^2}\bigl(g(z_{m,h})^2+o(1)\bigr).
\]
Put $y_m=(h-m\mu)/b_h$. Then $\lvert z_{m,h}-y_m\rvert\le M\lvert
b_h/a_m-1\rvert\to0$, so $g(z_{m,h})^2=g(y_m)^2+o(1)$. The $y_m$ with
$m\in\mathcal R_3$ are the points of an arithmetic progression with step
$\mu/b_h$ that lie in $[-M,M]$. Hence
\[
b_h\Sigma_3=2\bigl(1+o(1)\bigr)\frac\mu{b_h}
\sum_{m\in\mathcal R_3}\bigl(g(y_m)^2+o(1)\bigr)\to2\int_{-M}^Mg^2 .
\]
Adding the 4 regions,
\[
\limsup_{h\to\infty}\Bigl\lvert b_hC_n-2\int g^2\Bigr\rvert
\le2\int_{\lvert z\rvert>M}g^2+18c_0\bigl(\Prob(G\ge M-1)+\Prob(G\le-M)\bigr),
\]
and letting $M\to\infty$ gives $C_n\sim2b_h^{-1}\int g^2$.
\qed

\subsection{Proof of Theorem~\texorpdfstring{\ref{thm:alpha2}}{4.6}}

Let $g$ be the standard normal density, so that $\int g^2=1/(2\sqrt\pi)$. The
proof of Lemma~\ref{lem:C-firstpassage} used only $\delta_m\to0$,
$c_0<\infty$, the uniform continuity and decay of $g$, $T_J\to0$,
$a_{m-1}\to\infty$ and $a_m/a_{m-1}\to1$. So that lemma holds for $\alpha\ge2$ with the norming
below. With Lemma~\ref{lem:C-bigjump}(b), the proof of
Theorem~\ref{thm:center} then gives $C_n\sim1/(\sqrt\pi\,b_h)$, once we check
its other facts.

(a) Let $\alpha>2$. Summation by parts gives
$\E L^2=\sum_k(2k-1)\bar F(k)=2\zeta(\alpha-1)-\zeta(\alpha)$, so $v$ is as
stated and finite. With $a(x)=\sqrt{vx}$ the central limit theorem gives
$(\tau_m-m\mu)/a_m\Rightarrow\mathcal N(0,1)$, and Gnedenko's theorem
\cite[Theorem~4.2.1]{ibragimovlinnik1971} gives $\delta_m\to0$ and
$c_0<\infty$. So $C_n\sim(\pi vh/\mu)^{-1/2}=\sqrt{2\mu/(\pi vn)}$. For the
stationary start Proposition~\ref{prop:edge} gives
$E_n=n^{1-\alpha}(1+O(1/n))/(2(\alpha-1)\mu)$, and dividing gives
$R_n\sim C_2(\alpha)n^{3/2-\alpha}$.

(b) Let $\alpha=2$, so that $\mu=\zeta(2)$ and $\E L^2=\infty$. The truncated
second moment $U(x)=\E[L^2;L\le x]=\sum_{k\le x}(2k+1)/(k+1)^2=2\log x+O(1)$
is slowly varying. So $L$ is in the domain of attraction of the normal law,
and $(\tau_m-m\mu)/a_m\Rightarrow\mathcal N(0,1)$ for any $a_m$ with
$mU(a_m)/a_m^2\to1$ \cite{ibragimovlinnik1971}. Centering the truncated
moment at $\mu$ changes it only by $O(1)$. For $a(x)=\sqrt{x\log x}$ we have
$mU(a(m))/a(m)^2=(\log m+\log\log m+O(1))/\log m\to1$. This $a$ is
increasing on $[3,\infty)$ and regularly varying of index $\frac12$, and we
take $a_m=a(m)$ for $m\ge3$ and $a_1=a_2=1$. Gnedenko's theorem again gives
$\delta_m\to0$ and $c_0<\infty$. So $C_n\sim1/(\sqrt\pi\,b_h)$ with
$b_h^2=(h/\mu)\log(h/\mu)\sim n\log n/(2\mu)$, that is,
$C_n\sim\sqrt{2\mu/(\pi n\log n)}$. For the stationary start
$E_n=n^{-1}(1+O(1/n))/(2\mu)$ by Proposition~\ref{prop:edge}, so
$R_n\sim\sqrt{\pi\log n/(8\mu^3n)}$.

In both cases a fresh start gives the smaller edge $E_n=\frac12n^{-\alpha}$, so
$R_n\to0$ for either start.
\qed

\subsection{Proof of Theorem~\texorpdfstring{\ref{thm:lamperti}}{4.7}}

Let $0<\alpha<1$ with a fresh start. Then $u^D_m=u_m$ and $q^D_m=q_m$, so
Proposition~\ref{prop:tworenewal} gives
$C_n=\sum_{m\ge1}u_m(h)\bigl[q_m(h)+q_{m+1}(h)\bigr]$. Here
$a_m=m^{1/\alpha}$, $g$ is the density of Lemma~\ref{lem:stable}(v), and
$\delta_m\to0$. Since $g$ is bounded and $g(y)=O(y^{-1-\alpha})$, we have
$\int_0^\infty yg(y)^2\,dy<\infty$. Put $t_m=h/a_m$ and fix $M>1$.

\emph{Step 1.} Let $g_1(t)=\int_0^tw^{-\alpha}g(t-w)\,dw$. We claim that
$mq_m(h)=g_1(t_m)+o(1)$ as $h\to\infty$, uniformly in $m$ with
$t_m\in[1/(2M),2M]$. Such $m$ tend to infinity. Since
$\lvert k/a_{m-1}-k/a_m\rvert\le2M(a_m/a_{m-1}-1)\to0$ for $k\le h$, the
local limit theorem for $u_{m-1}$ and the uniform continuity of $g$ give
$a_mu_{m-1}(k)=g(k/a_m)+o(1)$ uniformly in $0\le k\le h$. As
$m\,j^{-\alpha}=(j/a_m)^{-\alpha}$ and $u_{m-1}(0)=0$, the sum for $q_m$ in
\eqref{eq:C-conv} becomes
\[
mq_m(h)=\frac1{a_m}\sum_{j=1}^{h-1}\Bigl(\frac j{a_m}\Bigr)^{-\alpha}
\Bigl(g\bigl(t_m-\tfrac j{a_m}\bigr)+o(1)\Bigr).
\]
For $x>0$ we have
$\frac1{a_m}\sum_{j\le xa_m}(j/a_m)^{-\alpha}\le x^{1-\alpha}/(1-\alpha)$.
With $x=t_m\le2M$ this shows that the $o(1)$ terms add up to $o(1)$. With
$x=\eta$ it shows that the terms with $j\le\eta a_m$ are $O(\eta^{1-\alpha})$,
and so is $\int_0^\eta w^{-\alpha}g(t_m-w)\,dw$. The other terms form a
Riemann sum with mesh $1/a_m$ of a function that is uniformly continuous on
$\{\eta\le w\le t\le2M\}$. Letting $\eta\to0$ proves the claim. Next, the
Laplace transform of $g_1$ is
$\Gamma(1-\alpha)s^{\alpha-1}e^{-\Gamma(1-\alpha)s^\alpha}$. This is
$\frac1\alpha$ times $-\frac d{ds}e^{-\Gamma(1-\alpha)s^\alpha}$, the
transform of $tg(t)$. Both functions are continuous on $(0,\infty)$, so
$g_1(t)=tg(t)/\alpha$.

\emph{Step 2.} If $t_m>M$, then $h>Ma_m$ and Lemma~\ref{lem:C-bigjump}(c)
gives $u_m(h)\le C_Bmh^{-\alpha-1}\le C_BM^{-\alpha}/h$. Since
$\sum_m(q_m(h)+q_{m+1}(h))\le2$, these $m$ contribute at most
$2C_BM^{-\alpha}/h$ to $C_n$. If $t_m<1/M$, then $u_m(h)\le c_0/a_m<c_0/(Mh)$,
and these $m$ contribute at most $2c_0/(Mh)$. For the remaining $m$ we have
$t_{m+1}=t_m(1+o(1))\in[1/(2M),2M]$ and $(m+1)/m\to1$, and $g_1$ is uniformly
continuous on $[1/(2M),2M]$. So Step 1, $1/a_m=t_m/h$ and $\delta_m\to0$ give,
uniformly,
\[
u_m(h)\bigl[q_m(h)+q_{m+1}(h)\bigr]
=\frac{t_m}h\bigl(g(t_m)+o(1)\bigr)\,\frac2m\Bigl(\frac{t_mg(t_m)}\alpha+o(1)\Bigr)
=\frac2{\alpha hm}\bigl(t_m^2g(t_m)^2+o(1)\bigr).
\]
As $m$ runs over these values, $\log m=\alpha\log(h/t_m)$ runs over an
interval of length $2\alpha\log M+o(1)$ with steps $(1+o(1))/m$. So the sum of
$1/m$ is $O(\log M)$, the $o(1)$ terms contribute $o(1/h)$, and the rest is a
Riemann sum in the variable $\log m$. It tends to
$\frac2{\alpha h}\int_{1/M}^Mt^2g(t)^2\,\frac{\alpha\,dt}t$. Hence
\[
hC_n=2\int_{1/M}^Mtg(t)^2\,dt+O(M^{-\alpha}+M^{-1})+o(1),
\]
and letting $M\to\infty$ gives $nC_n=2hC_n\to4\int_0^\infty tg(t)^2\,dt$.

\emph{Step 3.} Let $G_1,G_2$ be independent copies of $G$ and
$Y=G_1/(G_1+G_2)$. The map $(G_1,G_2)\mapsto(Y,G_1+G_2)$ has inverse
$(y,r)\mapsto(yr,(1-y)r)$ with Jacobian $r$. So $Y$ has the density
$f(y)=\int_0^\infty r\,g(yr)\,g((1-y)r)\,dr$ on $(0,1)$, and
$f(\frac12)=\int_0^\infty r\,g(r/2)^2\,dr=4\int_0^\infty tg(t)^2\,dt$. Hence
$nC_n\to f(\frac12)$. Note that $f$ is continuous at $\frac12$. Indeed, for
$y\in[\frac14,\frac34]$ both $yr$ and $(1-y)r$ are at least $r/4$, so by the
tail in Lemma~\ref{lem:stable}(v) the integrand is at most
$r\min(\lVert g\rVert_\infty^2,Cr^{-2-2\alpha})$. This is integrable, and
dominated convergence applies.

We now identify the law of $Y$. It does not change if $G$ is multiplied by a
constant, so we may take $\E e^{-sG}=e^{-s^\alpha}$. For $x>0$ and $0<p<1$,
\[
x^p=\frac p{\Gamma(1-p)}\int_0^\infty s^{-p-1}(1-e^{-sx})\,ds,\qquad
x^{-p}=\frac1{\Gamma(p)}\int_0^\infty s^{p-1}e^{-sx}\,ds .
\]
Taking expectations and substituting $w=s^\alpha$ gives
$\E G^p=\Gamma(1-p/\alpha)/\Gamma(1-p)$ for $\lvert p\rvert<\alpha$ (see also
\cite{zolotarev1986}). Let $W=G_1/G_2$, so that $Y=W/(1+W)$. By independence
and $\Gamma(1-z)\Gamma(1+z)=\pi z/\sin(\pi z)$,
\[
\E W^p=\frac{\Gamma(1-\frac p\alpha)\Gamma(1+\frac p\alpha)}{\Gamma(1-p)\Gamma(1+p)}
=\frac{\sin\pi p}{\alpha\sin(\pi p/\alpha)},\qquad\lvert p\rvert<\alpha .
\]
Let
$\rho_\alpha(w)=\frac{\sin\pi\alpha}\pi\,\frac{w^{\alpha-1}}{w^{2\alpha}+2w^\alpha\cos\pi\alpha+1}$
for $w>0$. We use the classical integral
\[
\int_0^\infty\frac{v^\beta\,dv}{v^2+2v\cos\theta+1}
=\frac{\pi\sin\beta\theta}{\sin\theta\,\sin\pi\beta},
\qquad\lvert\beta\rvert<1,\ 0<\theta<\pi,
\]
read as $\theta/\sin\theta$ at $\beta=0$. For $-1<\beta<0$ it follows from
partial fractions and
$\int_0^\infty v^\beta(v+e^{\pm i\theta})^{-1}dv=-\pi e^{\pm i\beta\theta}/\sin(\pi\beta)$,
and the substitution $v\mapsto1/v$ shows that both sides are even in
$\beta$. With $v=w^\alpha$, $\beta=p/\alpha$ and $\theta=\pi\alpha$ it gives
$\int_0^\infty w^p\rho_\alpha(w)\,dw=\E W^p$ for $\lvert p\rvert<\alpha$. At
$p=0$ this shows that $\rho_\alpha$ is a probability density. So $\log W$ and
the logarithm of a variable with density $\rho_\alpha$ have the same moment
generating function near $0$, and $W$ has density $\rho_\alpha$. Then $Y$ has
density $\rho_\alpha(\frac y{1-y})(1-y)^{-2}=\ell_\alpha(y)$. Both $f$ and
$\ell_\alpha$ are densities of $Y$ that are continuous at $\frac12$, so
\[
f(\tfrac12)=\ell_\alpha(\tfrac12)=\frac{2\sin\pi\alpha}{\pi(1+\cos\pi\alpha)}
=\frac2\pi\tan\frac{\pi\alpha}2 .
\]
Dividing $E_n=\frac12n^{-\alpha}$ by $C_n$ gives the formula for $R_n$.

\emph{Step 4.} Put $c=\cos\frac{\pi\alpha}2>0$, so that $1+\cos\pi\alpha=2c^2$
and $\sin(\pi\alpha)/(\pi c^2)=\frac2\pi\tan\frac{\pi\alpha}2$. For
$y=\frac12+\eta$ the binomial series give
$(y(1-y))^{\alpha-1}=4^{1-\alpha}\bigl(1-4(\alpha-1)\eta^2+O(\eta^4)\bigr)$ and
\[
y^{2\alpha}+2y^\alpha(1-y)^\alpha\cos\pi\alpha+(1-y)^{2\alpha}
=4^{1-\alpha}c^2\Bigl(1+\frac{4\alpha(\alpha-c^2)}{c^2}\,\eta^2+O(\eta^4)\Bigr).
\]
Hence
\[
\ell_\alpha(\tfrac12+\eta)=\frac2\pi\tan\frac{\pi\alpha}2
\Bigl[1+\frac{4(\cos^2\frac{\pi\alpha}2-\alpha^2)}{\cos^2\frac{\pi\alpha}2}\,\eta^2+O(\eta^4)\Bigr].
\]
The function $\alpha\mapsto\cos\frac{\pi\alpha}2-\alpha$ decreases strictly
on $(0,1)$ from $1$ to $-1$ and vanishes at $\alpha_c$. So the coefficient of
$\eta^2$ is positive for $\alpha<\alpha_c$ and negative for $\alpha>\alpha_c$,
and $y=\frac12$ is a strict local minimum of $\ell_\alpha$ in the first case
and a strict local maximum in the second. This agrees with
\cite[eq.~(7.6)]{godrecheluck2001}.
\qed
